\documentclass[12pt]{article}
\usepackage[utf8]{inputenc}
\usepackage[T1]{fontenc}
\usepackage{hyperref}
\title{Zoque, Tabasco Language Resources for K-12 Curriculum Development}
\author{Compiled by Claude Sonnet 5.5 (claude-sonnet-5-5) for Andr\'as Kornai}
\date{}
\begin{document}
\maketitle

\section{Introduction}
Tabasco Zoque (ISO 639-3: zoq; Glottocode taba1264) is a Mixe-Zoquean language of Tabasco, Mexico. It is also known as Ayapa Zoque or Ayapaneco. It has only a handful of speakers; UNESCO's atlas reports 24 speakers, and the survey figure is about 2. It is distinct from the Popoluca varieties, which are separate languages.

\section{Primary Educational Resources}
No primers, textbooks or school programmes were found. This section is empty.

\section{Linguistic and Reference Materials}
The Archive of the Indigenous Languages of Latin America (AILLA) holds a Zoque, Ayapaneco collection. A search summary described it as 92 audio recordings, 169 PDF documents and 25 text files, and I did not check these counts. \url{https://ailla-legacy.lib.utexas.edu/islandora/object/ailla%3A124491}

The UNESCO Atlas of World Languages entry lists digital audio and video materials, but few descriptions of the language. No published grammar or dictionary was confirmed. \url{https://ar.wal.unesco.org/en/countries/mexico/languages/tabasco-zoque}

\section{Religious and Community Resources}
None found.

\section{Digital and Audio Resources}
See the AILLA and UNESCO entries above. The Wikipedia page on the Zoque languages gives the family context. \url{https://en.wikipedia.org/wiki/Zoque_languages}

\section{Curriculum Development Context}
The language is critically endangered and thought to be used only by the elderly. It is not known to be taught in schools. With so few speakers, any programme would mean documentation and revitalisation work built on the archived recordings.

\end{document}